% !TeX root = ../../ONMIPA Alin.tex
% Acuan: Zhang (2009), Bab 1, khususnya Soal 1.1--1.7, 1.13--1.14,
% 1.22, 1.38, 1.41, 1.43, 1.44, dan 1.51; Zhang (2011), bagian 1.1.
% Pembahasan dibatasi sampai subruang; tidak membahas basis dan dimensi.
\chapter{Ruang Linear}
\label{ch:ruang-linear}

\textbf{Ruang vektor} atau bisa disebut \textbf{ruang linear}  merupakan himpunan yang dilengkapi dengan dua operasi, yaitu penjumlahan dan perkalian skalar. Menurut penulis, materi ini cukup penting dan bisa dipelajari untuk awal-awal karena tidak terlalu me-\textit{reference} definisi-definisi lain. Materi ini juga menjadi dasar untuk memahami materi-materi selanjutnya, seperti kebebasan linear, basis, dimensi, dan transformasi linear.

\section{Definisi Ruang Vektor}
\label{linear:aksioma}
Anggota $V$ disebut \emph{vektor}, sedangkan anggota medan $\F$
disebut \emph{skalar}. Contoh medan adalah $\R$, $\C$, dan
bilangan rasional $\mathbb Q$. Setiap skalar tak nol dalam suatu
medan mempunyai invers perkalian.

\begin{definisi}[][Ruang linear atas suatu medan][linear:def-ruang]
    Himpunan tak kosong $V$ disebut \emph{ruang linear atas $\F$}
    jika dilengkapi operasi
    \[
        +:V\times V\longrightarrow V,
        \qquad
        \cdot:\F\times V\longrightarrow V,
    \]
    yang memenuhi aksioma berikut untuk setiap $u,v,w\in V$
    dan $\alpha,\beta\in\F$:
    \begin{enumerate}[label=(\roman*),nosep]
        \item $u+v=v+u$;
        \item $(u+v)+w=u+(v+w)$;
        \item terdapat $0_V\in V$ sehingga $v+0_V=v$;
        \item untuk setiap $v$ terdapat $-v\in V$ sehingga $v+(-v)=0_V$;
        \item $\alpha(u+v)=\alpha u+\alpha v$;
        \item $(\alpha+\beta)v=\alpha v+\beta v$;
        \item $(\alpha\beta)v=\alpha(\beta v)$;
        \item $1v=v$, dengan $1$ identitas perkalian dalam $\F$.
    \end{enumerate}
\end{definisi}

Penulisan kedua operasi sebagai pemetaan ke $V$ sudah mencakup
\emph{ketertutupan}: $u+v$ dan $\alpha v$ harus tetap berada di $V$.
Jika ketertutupan ditulis sebagai dua syarat tersendiri, daftar di
atas menjadi sepuluh syarat. Kedua cara penulisan itu ekuivalen.

Vektor nol $0_V$ dan skalar nol $0\in\F$ mempunyai peran berbeda.
Pada $\R^n$, vektor nol adalah $(0,\ldots,0)^T$; pada ruang matriks,
ia adalah matriks nol; pada ruang fungsi, ia adalah fungsi yang
bernilai nol di setiap titik. Setelah jenis objeknya jelas, keduanya
sering ditulis dengan simbol $0$.

Untuk $v_1,\ldots,v_r\in V$ dan $\alpha_1,\ldots,\alpha_r\in\F$,
ekspresi $\alpha_1v_1+\cdots+\alpha_rv_r$ disebut
\emph{kombinasi linear} dari vektor-vektor tersebut.
Penjumlahan ini mempunyai hingga banyak suku dan selalu
menghasilkan anggota $V$ karena ketertutupan kedua operasi.

\begin{exam}[Medan skalar harus disebutkan]
    Tentukan apakah $\C$ merupakan ruang linear atas $\C$ dan atas
    $\R$. Apakah $\R$ merupakan ruang linear atas $\C$ dengan
    penjumlahan dan perkalian biasa?
\end{exam}
\begin{solution}
    Himpunan $\C$ merupakan ruang linear atas kedua medan tersebut:
    penjumlahan kompleks dan perkalian oleh skalar yang diperbolehkan
    tetap menghasilkan bilangan kompleks, dan aksiomanya mengikuti
    sifat aritmetika kompleks.
    Namun, $\R$ tidak tertutup terhadap perkalian oleh skalar kompleks,
    sebab $i\cdot1=i\notin\R$. Dengan operasi yang disebutkan,
    $\R$ bukan ruang linear atas $\C$.
\end{solution}

Secara umum, ruang linear atas $\C$ juga dapat dipandang sebagai
ruang linear atas $\R$ dengan membatasi skalar yang digunakan.
Arah sebaliknya memerlukan definisi perkalian oleh skalar kompleks
dan tidak berlaku secara otomatis.

\section{Contoh \emph{Well-Known}}
\label{linear:contoh}

\subsection{Koordinat, matriks, polinom, fungsi, dan barisan}

\begin{enumerate}[label=(\alph*)]
    \item \textbf{Vektor koordinat.} Himpunan $\F^n$ merupakan
          ruang linear dengan penjumlahan dan perkalian skalar per
          koordinat. Untuk perkalian $Ax$, vektor $x$ ditulis sebagai kolom.
    \item \textbf{Matriks berukuran tetap.} Himpunan
          $\Mat_{m\times n}(\F)$ merupakan ruang linear atas $\F$.
          Ukuran harus tetap agar penjumlahan terdefinisi.
          Perkalian dua matriks bukan bagian dari aksioma ruang linear.
    \item \textbf{Polinom.} Himpunan $\F[t]$ berisi semua polinom
          dengan koefisien dalam $\F$. Untuk $n\ge0$, tuliskan
          \[
              \mathcal P_{\le n}(\F)
              =\{a_0+a_1t+\cdots+a_nt^n:a_j\in\F\}.
          \]
          Polinom nol ikut termuat. Penjumlahan dan perkalian skalar
          tidak menaikkan derajat melampaui $n$, sehingga himpunan
          ini juga merupakan ruang linear. Notasi $\le n$ menegaskan
          bahwa batas derajatnya termasuk $n$.
    \item \textbf{Fungsi.} Untuk himpunan tak kosong $S$,
          himpunan semua fungsi $S\to\F$, ditulis $\F^S$,
          mempunyai operasi titik demi titik:
          \[
              (f+g)(s)=f(s)+g(s),\qquad
              (\alpha f)(s)=\alpha f(s).
          \]
          Aksioma berlaku karena dapat diperiksa pada setiap $s\in S$.
          Himpunan $C([a,b],\R)$ dari fungsi real kontinu pada
          $[a,b]$, $a<b$, juga merupakan ruang linear, sebab jumlah
          dan kelipatan skalar fungsi kontinu tetap kontinu.
    \item \textbf{Barisan.} Himpunan $\F^{\mathbb N}$ dari semua
          barisan $(a_1,a_2,\ldots)$ merupakan ruang linear dengan
          operasi per suku. Tidak ada syarat kekonvergenan dalam
          definisi ruang ini.
\end{enumerate}

Pada contoh-contoh tersebut, objeknya berbeda tetapi mekanisme
pembuktiannya sama: operasi diwarisi dari aritmetika dalam $\F$,
lalu diterapkan pada setiap koordinat, entri, koefisien, atau nilai fungsi.

\subsection{Bilangan positif sebagai vektor}

\begin{exam}[Penjumlahan berupa perkalian]
    Pada $V=\R_{>0}$, definisikan
    \[
        u\oplus v=uv,\qquad \alpha\odot u=u^\alpha
        \quad(u,v>0,\ \alpha\in\R).
    \]
    Buktikan bahwa $V$ merupakan ruang linear atas $\R$.
    Tentukan vektor nol dan invers penjumlahan dari $u$.
\end{exam}
\begin{solution}
    Kedua operasi tertutup karena hasilnya tetap positif.
    Sifat komutatif dan asosiatif $\oplus$ mengikuti perkalian real.
    Identitas penjumlahannya adalah $1$, dan invers penjumlahan $u$
    adalah $u^{-1}$, sebab $u\oplus1=u$ dan $u\oplus u^{-1}=1$.
    Hukum pangkat memberi
    \begin{align*}
        \alpha\odot(u\oplus v)
         & =(uv)^\alpha=u^\alpha v^\alpha
        =(\alpha\odot u)\oplus(\alpha\odot v), \\
        (\alpha+\beta)\odot u
         & =u^{\alpha+\beta}=u^\alpha u^\beta
        =(\alpha\odot u)\oplus(\beta\odot u),  \\
        (\alpha\beta)\odot u
         & =(u^\beta)^\alpha
        =\alpha\odot(\beta\odot u),
    \end{align*}
    serta $1\odot u=u$. Jadi seluruh aksioma terpenuhi.
    Di ruang ini, $0\odot u=1$ memang benar: bilangan $1$
    adalah vektor nol untuk operasi $\oplus$.
\end{solution}

Fungsi $u\mapsto\log u$ membantu membaca struktur tadi:
$\log(u\oplus v)=\log u+\log v$ dan
$\log(\alpha\odot u)=\alpha\log u$.
Dengan operasi biasa, $\R_{>0}$ tidak menjadi ruang linear atas
$\R$, sebab perkalian oleh skalar nol atau negatif keluar dari himpunan.

\subsection{Vektor nol dapat berada di titik lain}

\begin{exam}[Menggeser operasi pada ruang koordinat]
    Tetapkan $a\in\F^n$. Pada himpunan $\F^n$, definisikan
    \[
        u\oplus v=u+v-a,\qquad
        \alpha\odot u=a+\alpha(u-a).
    \]
    Tunjukkan bahwa operasi ini membentuk ruang linear atas $\F$.
\end{exam}
\begin{solution}
    Pemetaan $T(u)=u-a$ bijektif dari $\F^n$ ke dirinya dan memenuhi
    \[
        T(u\oplus v)=T(u)+T(v),\qquad
        T(\alpha\odot u)=\alpha T(u).
    \]
    Maka setiap identitas aksioma, setelah dikenai $T$, menjadi
    identitas pada ruang koordinat biasa. Karena $T$ injektif,
    identitas sebelum dikenai $T$ juga berlaku.
    Vektor nolnya adalah $a$, sebab $u\oplus a=u$.
    Invers penjumlahan dari $u$ adalah $2a-u$.
\end{solution}

Kedua contoh tersebut menunjukkan bahwa himpunan saja belum
menentukan suatu ruang linear. Medan dan kedua operasi harus diketahui.

\section{Definisi dan Kriteria Subruang}
\label{linear:subruang}

\begin{definisi}[][Subruang linear][linear:def-subruang]
    Misalkan $V$ ruang linear atas $\F$. Himpunan $W\subseteq V$
    disebut \emph{subruang linear} dari $V$ jika $W$ sendiri
    merupakan ruang linear atas $\F$ dengan penjumlahan dan
    perkalian skalar yang diwarisi dari $V$.
\end{definisi}

Frasa ``diwarisi dari $V$'' penting: kita tidak boleh mengganti
operasi atau medan skalar ketika menguji apakah $W$ subruang.
Misalnya, ruang bilangan positif dengan operasi $\oplus,\odot$
di atas bukan subruang dari $\R$ yang menggunakan operasi biasa.

\begin{teorema}[][Uji subruang][linear:uji-subruang]
    Misalkan $W\subseteq V$. Pernyataan berikut ekuivalen:
    \begin{enumerate}[label=(\roman*),nosep]
        \item $W$ merupakan subruang dari $V$;
        \item $0_V\in W$, $u+v\in W$ untuk semua $u,v\in W$,
              dan $\alpha u\in W$ untuk semua $\alpha\in\F$, $u\in W$;
        \item $W\ne\varnothing$ dan
              $\alpha u+\beta v\in W$ untuk semua $u,v\in W$
              serta $\alpha,\beta\in\F$.
    \end{enumerate}
\end{teorema}
\begin{proof}
    Jika $W$ subruang, vektor nolnya sama dengan vektor nol $V$:
    identitas $w+0_W=w$ dalam $V$ dan pencoretan memberi $0_W=0_V$.
    Ketertutupan kedua operasi mengikuti definisi, sehingga (i)
    mengakibatkan (ii). Dari (ii), perkalian skalar dan penjumlahan
    berurutan memberi (iii).

    Sekarang anggap (iii). Ambil satu $w\in W$.
    Dengan $\alpha=\beta=0$ diperoleh $0_V\in W$.
    Pilihan $\alpha=\beta=1$ memberi ketertutupan penjumlahan,
    sedangkan $\beta=0$ memberi ketertutupan perkalian skalar.
    Khususnya, $-u=(-1)u\in W$ untuk setiap $u\in W$.
    Aksioma asosiatif, komutatif, distributif, dan identitas skalar
    sudah berlaku di $V$, sehingga juga berlaku pada anggota $W$.
    Dengan demikian $W$ merupakan ruang linear.
\end{proof}

Syarat tak kosong tidak boleh dihilangkan. Pernyataan
``untuk setiap $u,v\in W$'' otomatis benar apabila $W$ kosong,
tetapi himpunan kosong tidak mempunyai vektor nol.
Subruang $\{0_V\}$ dan $V$ selalu ada. Subruang $W$ disebut
\emph{sejati} jika $W\ne V$; dengan pengertian ini,
$\{0_V\}$ juga sejati apabila $V\ne\{0_V\}$.

\begin{prop}[Subruang di dalam subruang]
    Jika $U$ subruang dari $W$ dan $W$ subruang dari $V$,
    semuanya atas medan yang sama, maka $U$ subruang dari $V$.
\end{prop}
\begin{proof}
    Inklusi $U\subseteq W\subseteq V$ berlaku dan kedua operasi
    yang digunakan sama. Vektor nol berada di $U$, dan $U$
    tertutup terhadap kedua operasi; terapkan uji subruang.
\end{proof}

\begin{exam}[Garis homogen dan garis bergeser]
    Untuk $c\in\R$, tentukan kapan
    $W_c=\{(x,y)\in\R^2:2x-3y=c\}$ merupakan subruang.
\end{exam}
\begin{solution}
    Jika $c\ne0$, vektor $(0,0)$ tidak berada di $W_c$, sehingga
    $W_c$ bukan subruang.
    Jika $c=0$ dan $u=(x_1,y_1),v=(x_2,y_2)\in W_0$, maka
    \[
        2(\alpha x_1+\beta x_2)-3(\alpha y_1+\beta y_2)
        =\alpha(2x_1-3y_1)+\beta(2x_2-3y_2)=0.
    \]
    Karena $(0,0)\in W_0$, uji subruang berlaku.
    Jadi syaratnya tepat $c=0$.
\end{solution}

\begin{exam}[Dua jenis ketertutupan yang berbeda]
    Himpunan $\Z^2\subset\R^2$ tertutup terhadap penjumlahan,
    tetapi tidak terhadap semua skalar real: $(1,0)\in\Z^2$,
    sedangkan $\tfrac12(1,0)\notin\Z^2$.
    Sebaliknya,
    \[
        S=\{(x,y)\in\R^2:xy=0\}
    \]
    tertutup terhadap perkalian skalar, tetapi tidak terhadap penjumlahan:
    \[
        (1,0),(0,1)\in S,\qquad (1,0)+(0,1)\notin S.
    \]
    Kedua himpunan memuat nol, tetapi keduanya bukan subruang real.
\end{exam}

\section{Mengenali Subruang dari Syarat Pembentuknya}
\label{linear:kendala}

\subsection{Persamaan homogen pada vektor dan matriks}

\begin{prop}[Himpunan solusi sistem homogen]
    \label{linear:kernel}
    Untuk $A\in\Mat_{m\times n}(\F)$, himpunan
    \[
        W=\{x\in\F^n:Ax=0\}
    \]
    merupakan subruang dari $\F^n$.
    Lebih umum, untuk $A\in\Mat_{m\times n}(\F)$ dan
    $B\in\Mat_{p\times q}(\F)$ tetap,
    \[
        \{X\in\Mat_{n\times p}(\F):AXB=0\}
    \]
    merupakan subruang dari $\Mat_{n\times p}(\F)$.
\end{prop}
\begin{proof}
    Vektor atau matriks nol memenuhi persamaan masing-masing.
    Jika $Au=Av=0$, maka
    $A(\alpha u+\beta v)=\alpha Au+\beta Av=0$.
    Untuk versi matriks, jika $AXB=AYB=0$, maka
    \[
        A(\alpha X+\beta Y)B=\alpha AXB+\beta AYB=0.
    \]
    Kedua kesimpulan mengikuti uji subruang.
\end{proof}

Himpunan pada bagian pertama adalah $\ker A$ yang telah muncul
pada bab rank dan nulitas. Untuk sistem tak homogen $Ax=b$
dengan $b\ne0$, himpunan solusinya bukan subruang dari $\F^n$,
karena vektor nol tidak memenuhi persamaan. Jika tidak ada solusi,
himpunannya kosong dan juga bukan subruang.

\begin{exam}[Parameter dalam suatu subruang]
    Buktikan bahwa
    \[
        W=\{(x_1,x_2,x_3,x_4)\in\F^4:
        x_3=x_1+x_2,\ x_4=x_1-x_2\}
    \]
    merupakan subruang. Tentukan semua anggotanya dalam bentuk parameter.
\end{exam}
\begin{solution}
    Kedua syarat adalah persamaan linear homogen, sehingga
    \autoref{linear:kernel} dapat diterapkan.
    Dengan $s=x_1$ dan $t=x_2$, seluruh anggota $W$ berbentuk
    \[
        (s,t,s+t,s-t),\qquad s,t\in\F.
    \]
    Ini juga memberi pemeriksaan langsung: kombinasi dengan skalar
    $\alpha,\beta$ dari anggota berparameter $(s_1,t_1)$ dan
    $(s_2,t_2)$ mempunyai parameter
    $(\alpha s_1+\beta s_2,\alpha t_1+\beta t_2)$, sehingga tetap di $W$.
\end{solution}

\subsection{Kendala pada polinom dan fungsi}

Evaluasi di titik tetap menghormati kedua operasi:
$(\alpha p+\beta q)(a)=\alpha p(a)+\beta q(a)$.
Karena itu, syarat seperti $p(a)=0$, $p(a)=p(b)$, atau
$2p(0)-p(1)=0$ menghasilkan subruang polinom.
Demikian pula, jika turunan terdefinisi, sifat
$(\alpha p+\beta q)'=\alpha p'+\beta q'$ dapat digunakan.

\begin{exam}[Mengubah syarat polinom menjadi syarat koefisien]
    Dalam $\mathcal P_{\le3}(\R)$, tentukan seluruh polinom
    yang memenuhi $p(0)=p(1)$ dan $p'(0)=0$.
    Buktikan bahwa himpunannya merupakan subruang.
\end{exam}
\begin{solution}
    Tuliskan $p(t)=a+bt+ct^2+dt^3$.
    Syarat $p'(0)=0$ memberi $b=0$, sedangkan $p(0)=p(1)$
    memberi $b+c+d=0$. Jadi $c=-d$ dan
    \[
        p(t)=a+d(t^3-t^2),\qquad a,d\in\R.
    \]
    Polinom nol diperoleh saat $a=d=0$.
    Penjumlahan dan perkalian skalar pada polinom hanya menjumlahkan
    dan mengalikan kedua parameter, sehingga himpunan tersebut subruang.
\end{solution}

Berbeda dari derajat \emph{paling besar} $n$, syarat derajat
\emph{tepat} $n$ tidak membentuk subruang, sebab polinom nol
tidak memenuhi syarat itu. Syarat $p(0)=1$ juga langsung gagal
pada polinom nol. Sementara itu, polinom yang tak negatif pada
seluruh $\R$ memuat polinom nol, tetapi tidak tertutup terhadap
skalar negatif: polinom $1$ memenuhi syarat, sedangkan $-1$ tidak.

\begin{prop}[Fungsi genap dan fungsi ganjil]
    \label{linear:paritas}
    Di dalam ruang semua fungsi $\R\to\F$, masing-masing himpunan
    \begin{align*}
        E & =\{f:f(-t)=f(t)\text{ untuk setiap }t\in\R\}, \\
        O & =\{f:f(-t)=-f(t)\text{ untuk setiap }t\in\R\}
    \end{align*}
    merupakan subruang atas $\F$.
\end{prop}
\begin{proof}
    Fungsi nol termasuk kedua himpunan. Jika $f,g\in E$,
    maka untuk setiap $t$,
    \[
        (\alpha f+\beta g)(-t)
        =\alpha f(t)+\beta g(t)=(\alpha f+\beta g)(t).
    \]
    Jadi $\alpha f+\beta g\in E$. Jika $f,g\in O$,
    perhitungan yang sama memberi
    $(\alpha f+\beta g)(-t)=-(\alpha f+\beta g)(t)$.
\end{proof}

Syarat itu harus berlaku untuk \emph{setiap} $t$, bukan hanya
pada sejumlah titik yang dipilih. Jika fungsi juga dipersyaratkan
kontinu, pembuktian tetap berlaku karena kombinasi linear hingga
dari fungsi kontinu tetap kontinu.

\subsection{Ketergantungan pada medan skalar}

Untuk matriks simetris, transpose memenuhi
$(\alpha A+\beta B)^T=\alpha A^T+\beta B^T$.
Jadi matriks simetris dalam $\Mat_n(\F)$ membentuk subruang atas
$\F$, baik ketika $\F=\R$ maupun ketika $\F=\C$.

Untuk transpose konjugat, skalar ikut dikonjugatkan:
\[
    (\alpha A+\beta B)^*=
    \overline\alpha A^*+\overline\beta B^*.
\]
Akibatnya, matriks Hermitian membentuk subruang \emph{real}
dari $\Mat_n(\C)$ yang dipandang sebagai ruang atas $\R$.
Namun, himpunan itu bukan subruang kompleks: $I$ Hermitian,
sedangkan $(iI)^*=-iI\ne iI$.

\begin{exam}[Subruang kompleks yang seluruh anggotanya Hermitian]
    Misalkan $W$ subruang kompleks dari $\Mat_n(\C)$ dan setiap
    $A\in W$ Hermitian. Buktikan $W=\{0\}$.
\end{exam}
\begin{solution}
    Untuk setiap $A\in W$, ketertutupan terhadap skalar kompleks
    memberi $iA\in W$. Karena $A$ dan $iA$ keduanya Hermitian,
    \[
        iA=(iA)^*=-iA^*=-iA.
    \]
    Jadi $2iA=0$ dan $A=0$. Karena $0\in W$, diperoleh $W=\{0\}$.
\end{solution}

\subsection{Syarat nonlinear perlu diperiksa langsung}

Tidak semua syarat yang tampak homogen menghasilkan subruang.
Contohnya, matriks $2\times2$ berank paling besar satu tertutup
terhadap perkalian skalar, tetapi
\[
    \rank\begin{pmatrix}1&0\\0&0\end{pmatrix}
    =\rank\begin{pmatrix}0&0\\0&1\end{pmatrix}=1,
    \qquad \rank I_2=2.
\]
Himpunan tersebut tidak tertutup terhadap penjumlahan.
Sebaliknya, suatu persamaan nonlinear dapat saja mendefinisikan
subruang: dalam $\R^2$, syarat $x^2+y^2=0$ hanya dipenuhi oleh
$(0,0)$, sehingga himpunan solusinya merupakan subruang.
Yang harus diperiksa adalah himpunan solusinya, bukan sekadar
bentuk persamaan yang dipakai untuk menuliskannya.

\section{Irisan dan Gabungan Subruang}
\label{linear:irisan-gabungan}

\begin{teorema}[][Irisan sebarang keluarga subruang][linear:irisan]
    Jika setiap $W_j$, $j\in J$, merupakan subruang dari $V$,
    maka $\bigcap_{j\in J}W_j$ merupakan subruang dari $V$.
\end{teorema}
\begin{proof}
    Vektor nol berada di setiap $W_j$, sehingga berada di irisan.
    Jika $u,v$ berada di irisan, maka keduanya berada di setiap
    $W_j$. Ketertutupan setiap $W_j$ memberi
    $\alpha u+\beta v\in W_j$ untuk setiap $j$, sehingga kombinasi
    itu berada di irisan. Untuk keluarga kosong, irisan relatif
    terhadap semesta $V$ diartikan sebagai $V$ sendiri.
\end{proof}

Sebagai contoh, polinom yang memenuhi $p(0)=p(1)=p(2)=0$
merupakan irisan tiga subruang yang masing-masing ditentukan
oleh satu syarat evaluasi. Hasil ini juga berlaku untuk tak hingga
banyak syarat: irisan himpunan $\{p\in\R[t]:p(a)=0\}$ atas
semua $a\in\R$ adalah subruang yang hanya berisi polinom nol.
Untuk fungsi pada \autoref{linear:paritas}, berlaku $E\cap O=\{0\}$,
sebab fungsi yang sekaligus genap dan ganjil memenuhi $f(t)=-f(t)$.

\begin{teorema}[][Gabungan dua subruang][linear:gabungan-dua]
    Jika $U,W$ subruang dari $V$, maka
    \[
        U\cup W\text{ subruang}
        \iff U\subseteq W\ \text{atau}\ W\subseteq U.
    \]
\end{teorema}
\begin{proof}
    Jika salah satu termuat di dalam yang lain, gabungannya adalah
    subruang yang lebih besar.
    Sebaliknya, andaikan tidak ada inklusi di antara keduanya.
    Pilih $u\in U\setminus W$ dan $w\in W\setminus U$.
    Jika gabungannya subruang, $u+w$ harus berada di $U\cup W$.
    Jika $u+w\in U$, maka $w=(u+w)-u\in U$, kontradiksi.
    Jika $u+w\in W$, maka $u=(u+w)-w\in W$, juga kontradiksi.
\end{proof}

Karena itu, gabungan dua garis berbeda yang melalui titik asal
di $\R^2$ bukan subruang. Gabungannya memuat dua arah, tetapi
tidak memuat seluruh hasil penjumlahan vektor dari arah-arah tersebut.

\begin{prop}[Gabungan subruang yang meningkat]
    \label{linear:rantai}
    Jika $W_1\subseteq W_2\subseteq\cdots$ adalah barisan
    subruang dari $V$, maka $\bigcup_{j\ge1}W_j$ merupakan subruang.
\end{prop}
\begin{proof}
    Vektor nol berada di gabungan. Jika $u\in W_r$ dan $v\in W_s$,
    ambil $k=\max(r,s)$. Kedua vektor berada di $W_k$, sehingga
    $\alpha u+\beta v\in W_k$ dan karenanya berada di gabungan.
\end{proof}

Contohnya,
$\R[t]=\bigcup_{n\ge0}\mathcal P_{\le n}(\R)$.
Setiap polinom mempunyai derajat berhingga, tetapi tidak ada satu
batas derajat yang berlaku untuk seluruh anggota $\R[t]$.

\subsection{Gabungan hingga dan medan tak hingga}

\begin{teorema}[][Gabungan hingga subruang sejati][linear:penutup-hingga]
    Misalkan $V$ ruang linear atas medan tak hingga, khususnya
    $\R$ atau $\C$. Jika $W_1,\ldots,W_m$ subruang sejati dari $V$,
    maka
    \[
        V\ne\bigcup_{j=1}^m W_j.
    \]
\end{teorema}
\begin{proof}
    Gunakan induksi pada $m$. Untuk $m=1$, pernyataan mengikuti
    definisi subruang sejati.
    Untuk langkah induksi, pilih
    $x\notin W_1\cup\cdots\cup W_{m-1}$.
    Jika $x\notin W_m$, pembuktian selesai.
    Jika $x\in W_m$, pilih $y\notin W_m$ dan pertimbangkan
    vektor $y+tx$, $t\in\F$.

    Tidak satu pun vektor itu berada di $W_m$: jika $y+tx\in W_m$,
    maka $y=(y+tx)-tx\in W_m$, kontradiksi.
    Untuk setiap $j<m$, paling banyak ada satu nilai $t$ yang
    membuat $y+tx\in W_j$. Sebab, jika dua nilai berbeda $s,t$
    memenuhi syarat itu, maka $(s-t)x\in W_j$, sehingga
    $x\in W_j$, kontradiksi.
    Hanya ada hingga banyak nilai parameter yang terlarang.
    Karena $\F$ tak hingga, kita dapat memilih $t$ lain dan
    memperoleh vektor di luar seluruh $W_j$.
\end{proof}

\begin{prop}[Satu subruang harus memuat semuanya]
    \label{linear:termuat-gabungan}
    Di atas $\R$ atau $\C$, jika subruang $U$ memenuhi
    $U\subseteq W_1\cup\cdots\cup W_m$, dengan semua $W_j$
    subruang, maka $U\subseteq W_j$ untuk sedikitnya satu $j$.
    Akibatnya, gabungan hingga subruang merupakan subruang tepat
    ketika salah satu anggotanya memuat semua anggota lainnya.
\end{prop}
\begin{proof}
    Tulis $U=\bigcup_{j=1}^m(U\cap W_j)$.
    Jika setiap $U\cap W_j$ sejati di dalam $U$, teorema sebelumnya
    memberi kontradiksi. Maka $U=U\cap W_j$ untuk suatu $j$.
    Untuk pernyataan terakhir, terapkan hasil ini pada subruang
    $U=\bigcup_j W_j$; arah sebaliknya langsung.
\end{proof}

Sifat hingga pada keluarga subruang penting. Semua garis yang
melalui titik asal di $\R^2$, jika digabungkan, menutupi seluruh
$\R^2$, walaupun masing-masing garis merupakan subruang sejati.
Pada penggunaan teorema ini, medan skalar juga harus diperiksa;
pembuktian pemilihan parameter tadi memakai ketersediaan tak
hingga banyak skalar.

\section{Latihan Pilihan}
\label{linear:latihan}

Nyatakan medan dan operasi yang digunakan sebelum menguji suatu
himpunan. Untuk jawaban negatif, berikan kegagalan aksioma atau
contoh konkret yang melanggar ketertutupan.

\begin{latihan}\label{linear:lat-operasi}
    Tetapkan $p,q\in\R^n$, $n\ge1$. Pada himpunan $\R^n$, definisikan
    \[
        u\oplus v=u+v+p,\qquad
        \alpha\odot u=\alpha u+(\alpha-1)q.
    \]
    Tentukan syarat perlu dan cukup pada $p,q$ agar kedua operasi
    itu membentuk ruang linear atas $\R$. Tentukan pula vektor nol
    dan invers penjumlahan setiap vektor.
\end{latihan}

\begin{latihan}\label{linear:lat-kompleks}
    Misalkan $V$ ruang linear kompleks, dan $W\subseteq V$
    merupakan subruang ketika skalar dibatasi pada $\R$.
    \begin{enumerate}[label=(\alph*),nosep]
        \item Buktikan bahwa $W$ subruang kompleks jika dan hanya
              jika $iw\in W$ untuk setiap $w\in W$.
        \item Tentukan apakah
              \[
                  \mathcal H=\left\{
                  \begin{pmatrix}z&w\\-\overline w&\overline z\end{pmatrix}
                  :z,w\in\C\right\}
              \]
              merupakan subruang real dan apakah merupakan subruang
              kompleks dari $\Mat_2(\C)$.
    \end{enumerate}
\end{latihan}

\begin{latihan}\label{linear:lat-kuadratik}
    Untuk $\lambda\in\R$, definisikan
    \[
        W_\lambda=\{(x,y)\in\R^2:x^2+\lambda y^2=0\}.
    \]
    Tentukan semua $\lambda$ yang membuat $W_\lambda$ subruang real.
    Kemudian jawab pertanyaan yang sama untuk
    $\widetilde W_\lambda=\{(x,y)\in\C^2:x^2+\lambda y^2=0\}$
    sebagai subruang kompleks, dengan parameter $\lambda$ tetap real.
\end{latihan}

\begin{latihan}\label{linear:lat-polinom}
    Di dalam $\mathcal P_{\le4}(\R)$, pertimbangkan
    \begin{align*}
        U & =\{p:p(0)=p(1)=p(2)\}, \\
        W & =\{p:p(0)p(1)=0\}.
    \end{align*}
    Tentukan himpunan mana yang merupakan subruang.
    Tuliskan semua polinom dalam $U$ dalam bentuk parameter,
    tanpa menghitung basis atau dimensi.
\end{latihan}

\begin{latihan}\label{linear:lat-komutatif}
    Misalkan $N=\begin{pmatrix}0&1\\0&0\end{pmatrix}$ dan
    \[
        U=\{X\in\Mat_2(\R):XN=NX\},\qquad
        W=\{X\in\Mat_2(\R):XN=-NX\}.
    \]
    Buktikan keduanya merupakan subruang dan tentukan semua
    anggotanya secara eksplisit. Apakah $W$ tertutup terhadap
    perkalian dua matriks? Jelaskan kaitannya dengan definisi subruang.
\end{latihan}

\begin{latihan}\label{linear:lat-trace-kuadrat}
    Tentukan apakah masing-masing himpunan berikut merupakan
    subruang dari $\Mat_2(\R)$:
    \[
        S=\{X:\tr(X^2)=0\},\qquad
        T=\{X:X^T=X,\ \tr(X^2)=0\}.
    \]
    Semua operasi menggunakan penjumlahan dan perkalian skalar biasa.
\end{latihan}

\begin{latihan}\label{linear:lat-rekurensi}
    Untuk $c\in\R$, misalkan $W_c$ terdiri atas semua barisan real
    $(a_n)_{n\ge1}$ yang memenuhi
    \[
        a_{n+2}-a_{n+1}-a_n=c\qquad(n\ge1).
    \]
    Tentukan kapan $W_c$ merupakan subruang dari $\R^{\mathbb N}$.
    Untuk $W_0$, buktikan bahwa setiap pilihan $a_1,a_2$ menentukan
    tepat satu barisan, dan bahwa syarat $a_1=a_2=0$ hanya
    menghasilkan barisan nol.
\end{latihan}

\begin{latihan}\label{linear:lat-afine}
    Misalkan $W\subseteq V$, dengan $V$ ruang linear real,
    memenuhi $0_V\in W$ serta
    \[
        \alpha u+(1-\alpha)v\in W
        \quad\text{untuk setiap }u,v\in W,\ \alpha\in\R.
    \]
    Buktikan $W$ subruang.
    Apakah masih cukup jika ketertutupan tersebut hanya diketahui
    untuk $\alpha=\tfrac12$? Apakah syarat $0_V\in W$ boleh dihilangkan?
\end{latihan}

\begin{latihan}\label{linear:lat-barisan}
    Di dalam $\R^{\mathbb N}$, definisikan
    \begin{align*}
        C_{00} & =\{(a_n):a_n=0\text{ untuk semua }n\text{ yang cukup besar}\},       \\
        E      & =\{(a_n):a_n\text{ konstan untuk semua }n\text{ yang cukup besar}\}.
    \end{align*}
    Buktikan bahwa keduanya subruang dan bahwa
    $C_{00}\subsetneq E\subsetneq\R^{\mathbb N}$.
    Nyatakan $C_{00}$ sebagai gabungan meningkat subruang-subruang
    yang masing-masing mengizinkan hanya hingga banyak suku awal tak nol.
\end{latihan}

\begin{latihan}\label{linear:lat-garis}
    Buktikan bahwa setiap subruang sejati tak nol dari $\R^2$
    berbentuk $\{tu:t\in\R\}$ untuk suatu $u\ne0$.
    Gunakan hasil ini untuk menjelaskan mengapa gabungan tak hingga
    subruang sejati dapat sama dengan seluruh ruang.
\end{latihan}

\begin{latihan}\label{linear:lat-hindari}
    Diberikan vektor-vektor tak nol $a_1,\ldots,a_m\in\R^n$,
    dengan $m,n\ge1$. Buktikan bahwa ada $x\in\R^n$ sehingga
    \[
        a_j^Tx\ne0\qquad\text{untuk semua }j=1,\ldots,m.
    \]
    Berikan dua pembuktian: satu menggunakan teorema gabungan hingga,
    dan satu lagi dengan memilih $x=(1,t,\ldots,t^{n-1})^T$ secara tepat.
\end{latihan}

\begin{latihan}\label{linear:lat-nol-bersama}
    Buktikan kedua pernyataan berikut.
    \begin{enumerate}[label=(\alph*)]
        \item Misalkan $W$ subruang dari $\R^n$ dan setiap $w\in W$
              mempunyai sedikitnya satu koordinat nol. Maka terdapat
              satu indeks $j$ yang sama sehingga $w_j=0$ untuk semua $w\in W$.
        \item Misalkan $U$ subruang dari $\R[t]$, dan
              $a_1,\ldots,a_m$ bilangan real tetap yang berbeda, $m\ge1$.
              Jika setiap $p\in U$ mempunyai akar di antara
              $a_1,\ldots,a_m$, maka ada satu $a_j$ yang menjadi
              akar semua polinom dalam $U$.
    \end{enumerate}
\end{latihan}

\ifpembahasan
    \section{Pembahasan Latihan}
    \label{linear:pembahasan}

    \subsection*{Latihan \ref{linear:lat-operasi}: dua cara menentukan vektor nol}
    Identitas untuk operasi $\oplus$ harus memenuhi $u\oplus e=u$,
    yakni $u+e+p=u$. Jadi $e=-p$.
    Jika kedua operasi membentuk ruang linear, identitas $0\odot u=e$
    harus berlaku. Namun, $0\odot u=-q$, sehingga $p=q$ merupakan
    syarat perlu.
    Jika $p=q$, ambil $a=-p$ pada contoh pergeseran operasi:
    \[
        u\oplus v=u+v-a,\qquad
        \alpha\odot u=a+\alpha(u-a).
    \]
    Contoh tersebut membuktikan kecukupan. Dengan demikian syaratnya
    tepat $p=q$; vektor nolnya $-p$ dan invers penjumlahan dari $u$
    adalah $-u-2p$.

    \subsection*{Latihan \ref{linear:lat-kompleks}: cukup menguji perkalian oleh $i$}
    Jika $W$ subruang kompleks, tentu $iw\in W$.
    Sebaliknya, tulis $\alpha=a+bi$ dengan $a,b\in\R$.
    Karena $W$ subruang real dan $iw\in W$, diperoleh
    \[
        \alpha w=aw+b(iw)\in W.
    \]
    Ketertutupan penjumlahan sudah diketahui, sehingga $W$ subruang kompleks.

    Untuk $\mathcal H$, matriks nol diperoleh dari $z=w=0$.
    Penjumlahan menjumlahkan parameter $z,w$, dan perkalian oleh
    skalar real mempertahankan hubungan konjugasi pada kedua entri
    baris bawah. Jadi $\mathcal H$ subruang real.
    Matriks $I_2$ berada di $\mathcal H$, tetapi $iI_2$ tidak:
    entri kiri atasnya memaksa $z=i$, sehingga entri kanan bawah
    seharusnya $\overline z=-i$, bukan $i$.
    Maka $\mathcal H$ bukan subruang kompleks.

    \subsection*{Latihan \ref{linear:lat-kuadratik}: medan mengubah himpunan solusi}
    Untuk skalar real, jika $\lambda>0$, persamaan hanya mempunyai
    solusi $(0,0)$. Jika $\lambda=0$, solusinya seluruh pasangan
    $(0,y)$; kedua himpunan merupakan subruang.
    Jika $\lambda<0$, letakkan $c=\sqrt{-\lambda}>0$.
    Vektor $(c,1)$ dan $(-c,1)$ memenuhi persamaan, tetapi jumlahnya
    $(0,2)$ tidak, sebab $0^2+\lambda\,2^2=4\lambda\ne0$.
    Jadi $W_\lambda$ subruang real tepat ketika $\lambda\ge0$.

    Di atas $\C$, untuk setiap $\lambda\ne0$ terdapat $c\ne0$
    dengan $c^2=-\lambda$. Kedua vektor yang sama kembali memberi
    kegagalan ketertutupan penjumlahan.
    Karena kasus $\lambda=0$ tetap menghasilkan $\{(0,y):y\in\C\}$,
    $\widetilde W_\lambda$ subruang kompleks tepat ketika $\lambda=0$.

    \subsection*{Latihan \ref{linear:lat-polinom}: tiga nilai yang sama dan satu perkalian}
    Syarat pada $U$ ekuivalen dengan dua persamaan linear homogen
    $p(0)-p(1)=0$ dan $p(1)-p(2)=0$, sehingga $U$ subruang.
    Jika $a=p(0)$, polinom $p(t)-a$ mempunyai akar $0,1,2$.
    Teorema faktor memberi
    \[
        p(t)=a+t(t-1)(t-2)(bt+c),\qquad a,b,c\in\R,
    \]
    karena derajat $p$ paling besar empat. Sebaliknya, setiap polinom
    berbentuk ini mempunyai nilai $a$ pada ketiga titik.

    Himpunan $W$ bukan subruang. Polinom $p(t)=t$ dan $q(t)=t-1$
    masing-masing berada di $W$, tetapi $p+q=2t-1$ mempunyai nilai
    $-1$ di $0$ dan $1$ di $1$, sehingga tidak berada di $W$.

    \subsection*{Latihan \ref{linear:lat-komutatif}: persamaan linear pada entri}
    Dengan $X=\begin{pmatrix}a&b\\c&d\end{pmatrix}$,
    \[
        XN=\begin{pmatrix}0&a\\0&c\end{pmatrix},\qquad
        NX=\begin{pmatrix}c&d\\0&0\end{pmatrix}.
    \]
    Maka
    \[
        U=\left\{\begin{pmatrix}a&b\\0&a\end{pmatrix}:a,b\in\R\right\},
        \qquad
        W=\left\{\begin{pmatrix}a&b\\0&-a\end{pmatrix}:a,b\in\R\right\}.
    \]
    Matriks nol berada di kedua himpunan, dan kombinasi linear
    mempertahankan bentuk masing-masing. Jadi keduanya subruang.
    Matriks $D=\diag(1,-1)$ berada di $W$, tetapi $D^2=I_2$ tidak
    berada di $W$. Ini tidak bertentangan dengan definisi subruang,
    karena definisi tersebut tidak mensyaratkan ketertutupan perkalian matriks.

    \subsection*{Latihan \ref{linear:lat-trace-kuadrat}: simetris memberi jumlah kuadrat}
    Ambil $A=\begin{pmatrix}0&1\\0&0\end{pmatrix}$ dan
    $B=\begin{pmatrix}0&0\\1&0\end{pmatrix}$.
    Karena $A^2=B^2=0$, keduanya berada di $S$.
    Namun, $(A+B)^2=I_2$, sehingga $\tr((A+B)^2)=2\ne0$.
    Jadi $S$ bukan subruang.
    Untuk matriks simetris real $X=\begin{pmatrix}a&b\\b&c\end{pmatrix}$,
    \[
        \tr(X^2)=a^2+2b^2+c^2.
    \]
    Nilainya nol tepat ketika $a=b=c=0$. Jadi $T=\{0\}$ merupakan subruang.

    \subsection*{Latihan \ref{linear:lat-rekurensi}: ketertutupan diperiksa pada setiap indeks}
    Jika $c\ne0$, barisan nol tidak memenuhi persamaan, sehingga
    $W_c$ bukan subruang. Untuk $c=0$, barisan nol termasuk $W_0$.
    Jika $(a_n),(b_n)\in W_0$ dan $d_n=\alpha a_n+\beta b_n$, maka
    \[
        d_{n+2}-d_{n+1}-d_n
        =\alpha(a_{n+2}-a_{n+1}-a_n)
        +\beta(b_{n+2}-b_{n+1}-b_n)=0
    \]
    untuk setiap $n$. Jadi $W_0$ subruang.
    Setelah $a_1,a_2$ dipilih, rumus $a_{n+2}=a_{n+1}+a_n$
    mendefinisikan setiap suku berikutnya secara tunggal.
    Induksi menunjukkan bahwa dua barisan dengan suku awal sama
    harus identik. Jika kedua suku awal nol, induksi juga memberi
    $a_n=0$ untuk setiap $n$.

    \subsection*{Latihan \ref{linear:lat-afine}: peran nol dan seluruh skalar}
    Dengan memilih $v=0_V$, syarat soal memberi $\alpha u\in W$
    untuk setiap skalar real $\alpha$.
    Untuk $u,v\in W$, pilihan $\alpha=\tfrac12$ memberi
    $(u+v)/2\in W$. Ketertutupan skalar yang baru dibuktikan
    kemudian memberi $u+v\in W$. Maka $W$ subruang.

    Jika hanya titik tengah yang diketahui berada di $W$, kesimpulan
    dapat gagal. Himpunan $\mathbb Q\subset\R$ memuat nol dan tertutup
    terhadap rata-rata dua anggotanya, tetapi
    $\sqrt2\cdot1\notin\mathbb Q$.
    Syarat memuat nol juga penting: garis
    $W=\{(x,1):x\in\R\}$ tertutup terhadap semua kombinasi
    $\alpha u+(1-\alpha)v$, tetapi bukan subruang karena tidak memuat $(0,0)$.

    \subsection*{Latihan \ref{linear:lat-barisan}: ambil batas indeks yang lebih besar}
    Jika dua barisan nol setelah indeks masing-masing $r$ dan $s$,
    kombinasi linearnya nol setelah $\max(r,s)$.
    Jika dua barisan konstan setelah indeks-indeks itu, dengan nilai
    akhir $a$ dan $b$, kombinasi linearnya konstan dengan nilai
    $\alpha a+\beta b$ setelah $\max(r,s)$.
    Kedua himpunan memuat barisan nol, sehingga keduanya subruang.
    Barisan yang akhirnya nol tentu akhirnya konstan, tetapi
    $(1,1,1,\ldots)\in E\setminus C_{00}$.
    Barisan $(1,2,3,\ldots)$ tidak berada di $E$, sehingga kedua
    inklusi pada soal bersifat sejati.
    Dengan
    \[
        V_k=\{(a_n):a_n=0\text{ untuk setiap }n>k\},\qquad k\ge1,
    \]
    setiap $V_k$ subruang, $V_k\subseteq V_{k+1}$, dan
    $C_{00}=\bigcup_{k\ge1}V_k$.

    \subsection*{Latihan \ref{linear:lat-garis}: dua arah berbeda memenuhi seluruh bidang}
    Misalkan $W$ subruang sejati tak nol. Pilih $u\in W\setminus\{0\}$.
    Ketertutupan skalar memberi $\{tu:t\in\R\}\subseteq W$.
    Andaikan ada $v\in W$ yang bukan kelipatan $u$.
    Matriks $[u\ v]$ berukuran dua mempunyai determinan tak nol:
    jika determinannya nol, hubungan antar kedua koordinat
    akan memaksa $v$ kelipatan $u$.
    Karena itu, untuk setiap $x\in\R^2$, sistem
    $[u\ v](\alpha,\beta)^T=x$ mempunyai solusi.
    Jadi setiap $x$ adalah $\alpha u+\beta v\in W$, bertentangan
    dengan $W$ sejati. Maka $W=\{tu:t\in\R\}$.
    Setiap vektor tak nol terletak pada garis yang dibentuk oleh
    kelipatan dirinya. Gabungan seluruh garis demikian, bersama
    vektor nol yang sudah termuat pada tiap garis, adalah $\R^2$.

    \subsection*{Latihan \ref{linear:lat-hindari}: hindari hingga banyak kendala}
    Untuk setiap $j$, definisikan $H_j=\{x:a_j^Tx=0\}$.
    Ini subruang sejati: karena $a_j\ne0$, terdapat koordinat tak nol,
    dan vektor satuan pada koordinat tersebut tidak berada di $H_j$.
    Teorema \ref{linear:penutup-hingga} menjamin ada
    $x\notin\bigcup_j H_j$, sehingga semua $a_j^Tx\ne0$.

    Untuk cara kedua, tulis $a_j=(a_{j1},\ldots,a_{jn})^T$ dan bentuk
    \[
        f_j(t)=a_{j1}+a_{j2}t+\cdots+a_{jn}t^{n-1}.
    \]
    Setiap $f_j$ polinom tak nol, sehingga hanya mempunyai hingga
    banyak akar real. Pilih $t$ di luar gabungan seluruh akar
    $f_1,\ldots,f_m$. Vektor $x=(1,t,\ldots,t^{n-1})^T$
    memenuhi $a_j^Tx=f_j(t)\ne0$ untuk setiap $j$.

    \subsection*{Latihan \ref{linear:lat-nol-bersama}: satu koordinat atau akar bersama}
    Untuk (a), definisikan subruang koordinat
    $H_j=\{x\in\R^n:x_j=0\}$.
    Hipotesis soal menyatakan $W\subseteq\bigcup_{j=1}^n H_j$.
    Proposisi \ref{linear:termuat-gabungan} memberi $W\subseteq H_j$
    untuk sedikitnya satu indeks $j$. Itulah koordinat yang selalu nol.

    Untuk (b), ambil $K_j=\{p\in\R[t]:p(a_j)=0\}$, yang merupakan
    subruang karena evaluasi menghormati penjumlahan dan perkalian skalar.
    Hipotesis menyatakan $U\subseteq\bigcup_{j=1}^m K_j$.
    Proposisi yang sama memberi $U\subseteq K_j$ untuk suatu $j$,
    sehingga $p(a_j)=0$ untuk semua $p\in U$.
    Perhatikan perubahan kuantornya: pilihan koordinat atau akar
    pada hipotesis boleh bergantung pada vektor; kesimpulannya memberi
    satu pilihan yang berlaku serentak untuk seluruh subruang.
\fi
